\documentclass[a4paper]{article}
\usepackage[top=1in, bottom=1.25in, left=1.25in, right=1.25in]{geometry}
\usepackage{amsmath}
\usepackage{multicol}
\usepackage{graphicx}
\usepackage[utf8]{inputenc}
\usepackage[english]{babel}
\setlength{\parskip}{0.03cm plus4mm minus3mm}
\RequirePackage{ltxcmds}[2010/12/07]

\usepackage{hyperref}
%opening
\title{MQAM transmitter}

\begin{document}

\maketitle
This block generates one or two optical signals. It can also generate a binary signal. It allows the user to controll the whole system. This block is a big block of code that combines all of the other in order to generate a optical signal. A schematic representation of this block is shown in figure \ref{MQAM_transmitter_block_diagram}.

\begin{figure}
	\centering
	\includegraphics[width=\textwidth]{MQAM_transmitter_block_diagram}
	\caption{Schematic representation of the block MQAM transmitter.}\label{MQAM_transmitter_block_diagram}
\end{figure}

\subsection*{Input parameters}

This block has a special set of functions that allow the user to change the basic configuration of the transmitter. This functions are listed bellow.

\begin{itemize}
	\item setMode(PseudoRandom);
	\item setBitPeriod(1.0/50e9);
	\linebreak (double)
	\item setPatternLength(3);
	\linebreak (int)
	\item setNumberOfBits(10000);
	\linebreak (long)
	\item setNumberOfSamplesPerSymbol(32);
	\linebreak (int)
	\item setRollOffFactor(0.9);
	\linebreak (double $\in$ [0,1])
	\item setIqAmplitudes(\{ \{ 1, 1 \}, \{ -1, 1 \}, \{ -1, -1 \}, \{ 1, -1 \} \});
	\item setOutputOpticalPower\_dBm(0);
	\item setSaveInternalSignals(true);
\end{itemize}


\subsection*{Methods}

MQamTransmitter(vector$<$Signal *$>$ \&inputSignal, vector$<$Signal *$>$ \&outputSignal);
\bigbreak

void set(int opt);
\bigbreak
void setMode(BinarySourceMode m) \{ B1.setMode(m); \};
\bigbreak
BinarySourceMode const getMode(void) \{ return B1.getMode(); \};
\bigbreak
void setProbabilityOfZero(double pZero) \{ B1.setProbabilityOfZero(pZero); \};
\bigbreak
double const getProbabilityOfZero(void) \{ return B1.getProbabilityOfZero(); \};
\bigbreak
void setBitStream(string bStream) \{ B1.setBitStream(bStream); \};
\bigbreak
string const getBitStream(void) \{ return B1.getBitStream(); \};
\bigbreak
void setNumberOfBits(long int nOfBits) \{ B1.setNumberOfBits(nOfBits); \}
\bigbreak
long int const getNumberOfBits(void) \{ return B1.getNumberOfBits();  \}
\bigbreak
void setPatternLength(int pLength) \{ B1.setPatternLength(pLength); \}
\bigbreak
int const getPatternLength(void) \{ return B1.getPatternLength(); \}
\bigbreak
void setBitPeriod(double bPeriod) \{B1.setBitPeriod(bPeriod);\};
double const getBitPeriod(void) \{ return B1.getBitPeriod(); \}
\bigbreak
void setM(int mValue)\{B2.m = mValue;\};
int const getM(void) \{ return B2.m; \};
\bigbreak
%void setIqAmplitudes(vector$<$t$_$iqValues$>$ iqAmplitudesValues) \{ B2.m = iqAmplitudesValues.size(); B2.iqAmplitudes.resize(B2.m);B2.iqAmplitudes = iqAmplitudesValues; \};
\bigbreak
%vector$<$t$_$iqValues$>$ const getIqAmplitudes(void) \{ return B2.iqAmplitudes; \};
\bigbreak
void setNumberOfSamplesPerSymbol(int n)\{ B3.setNumberOfSamplesPerSymbol(n); B4.setNumberOfSamplesPerSymbol(n); \};
\bigbreak
int const getNumberOfSamplesPerSymbol(void)\{ return B3.getNumberOfSamplesPerSymbol(); \};
\bigbreak
void setRollOffFactor(double rOffFactor)\{ B5.setRollOffFactor(rOffFactor); B6.setRollOffFactor(rOffFactor); \};
\bigbreak
double const getRollOffFactor(void)\{ return B5.getRollOffFactor(); \};
\bigbreak
void setSeeBeginningOfImpulseResponse(bool sBeginningOfImpulseResponse)\{ B5.setSeeBeginningOfImpulseResponse(sBeginningOfImpulseResponse); B6.setSeeBeginningOfImpulseResponse(sBeginningOfImpulseResponse); \};
\bigbreak
double const getSeeBeginningOfImpulseResponse(void)\{ return B5.getSeeBeginningOfImpulseResponse(); \};
\bigbreak
%void setOutputOpticalPower(t$_$real outOpticalPower) \{ B7.outputOpticalPower = outOpticalPower; \};
\bigbreak
%t$_$real const getOutputOpticalPower(void) \{ return B7.outputOpticalPower; \};
\bigbreak
%void setOutputOpticalPower$_$dBm(t$_$real outOpticalPower$_$dBm) \{ B7.outputOpticalPower = 1e-3*pow(10, outOpticalPower$_$dBm / 10); \};
\bigbreak
%t$_$real const getOutputOpticalPower$_$dBm(void) \{ return 10*log10(B7.outputOpticalPower/1e-3); \}

\subsection*{Functional description}

This block can generate one or two optical signals. If it generates two they have polarizations along perpendicular axis so that it is possible to combine them to produce a signal with any type of polarization. It also generates a binary signal that should be used to measure the Bit Error Rate (BER). 

\subsection*{Output Signals}

\subparagraph*{Number:} 1 or 2 optical and 1 binary

\subparagraph*{Type:} Optical signal

\subsection*{Example} 

\subsection*{Sugestions for future improvement}


\end{document}